\clearpage
\VerdierPhysicalPageAnchor{0246}%
\begingroup\small\setlength{\parskip}{0pt}\setlength{\emergencystretch}{1.5em}
\setlength{\abovedisplayskip}{4pt plus 2pt minus 2pt}%
\setlength{\belowdisplayskip}{4pt plus 2pt minus 2pt}%

\VerdierSemanticAnchor{ch4d-4-6-4}
\noindent\textbf{4.6.4.}\quad
Soit $X^{\bullet\bullet}$ un double complexe de $\mathcal A$
(chap.~I, 2.1.2). Le double complexe $X^{\bullet\bullet}$ définit un objet
de $\operatorname{comp}(\operatorname{comp}(\mathcal A))$ :

\VerdierSemanticAnchor{ch4d-4-6-4-double-complex}
\[
X^{\bullet\bullet}\colon\quad
\cdots\longrightarrow X^{\bullet,n-1}
\xrightarrow{d_2^{\bullet,n-1}}X^{\bullet,n}
\xrightarrow{d_2^{\bullet,n}}X^{\bullet,n+1}
\longrightarrow\cdots.
\]
(On utilise les notations de (4.1.9). Attention au signe !)

\VerdierSemanticAnchor{ch4d-4-6-4-b1}
Désignons par $B_1^p(X^{\bullet\bullet})$ le complexe simple :

\VerdierSemanticAnchor{ch4d-4-6-4-1}
\[
\left\{
\begin{aligned}
B_1^p(X^{\bullet\bullet})^i
  &=\operatorname{Im}(d_1^{p-1,i}),\\
d_{B_1^p(X^{\bullet\bullet})}^{i}
  &\text{ induit par }d_2^{p,i}.
\end{aligned}
\right.
\tag{4.6.4.1}
\]

\VerdierSemanticAnchor{ch4d-4-6-4-z1}
Désignons par $Z_1^p(X^{\bullet\bullet})$ le complexe simple :

\VerdierSemanticAnchor{ch4d-4-6-4-2}
\[
\left\{
\begin{aligned}
Z_1^p(X^{\bullet\bullet})^i
  &=\operatorname{Ker}(d_1^{p,i}),\\
d_{Z_1^p(X^{\bullet\bullet})}^{i}
  &\text{ induit par }d_2^{p,i}.
\end{aligned}
\right.
\tag{4.6.4.2}
\]

\VerdierSemanticAnchor{ch4d-4-6-4-h1}
Désignons par $H_1^p(X^{\bullet\bullet})$ le complexe simple :

\VerdierSemanticAnchor{ch4d-4-6-4-3}
\[
\left\{
\begin{aligned}
H_1^p(X^{\bullet\bullet})^i
  &=Z_1^p(X^{\bullet\bullet})^i/B_1^p(X^{\bullet\bullet})^i,\\
d_{H_1^p(X^{\bullet\bullet})}^{i}
  &\text{ induit par }d_2^{p,i}.
\end{aligned}
\right.
\tag{4.6.4.3}
\]

\VerdierSemanticAnchor{ch4d-4-6-4-simple-as-double}
Enfin, $Y$ étant un complexe simple de $\mathcal A$, nous désignons encore
par $Y$ le complexe double dont les composants sont nuls en deuxième degré
non nul et dont les composants et les premières différentielles sont ceux de
$Y$ en deuxième degré zéro.

\VerdierSemanticAnchor{ch4d-4-6-5}
\medskip\noindent\textbf{Définition 4.6.5.}\quad
Soit $Y$ un complexe simple de $\mathcal A$. Une \emph{résolution injective
de Cartan--Eilenberg} de $Y$ est un double complexe :

\VerdierSemanticAnchor{ch4d-4-6-5-double-complex}
\[
I^{\bullet\bullet}\colon\quad
\cdots\longrightarrow0\longrightarrow I^{\bullet,0}
\longrightarrow I^{\bullet,1}\longrightarrow I^{\bullet,2}
\longrightarrow\cdots,
\]
muni d'un morphisme appelé \emph{augmentation} :

\VerdierSemanticAnchor{ch4d-4-6-5-augmentation}
\[
\varepsilon\colon Y\longrightarrow I^{\bullet\bullet}
\]
tel que pour tout entier $p$, les complexes simples
$B_1^p(I^{\bullet\bullet})$ (4.6.4.1) et
$H_1^p(I^{\bullet\bullet})$ (4.6.4.3) munis des augmentations :

\VerdierSemanticAnchor{ch4d-4-6-5-augmentations}
\[
B_1^p(\varepsilon)\colon
\operatorname{Im}(d_Y^{p-1})\longrightarrow B_1^p(I^{\bullet\bullet}),
\]
\[
H_1^p(\varepsilon)\colon
H^p(Y)\longrightarrow H_1^p(I^{\bullet\bullet})
\]
\endgroup
